\chapter{Ethical analysis}
\label{app:ethics}

\noindent\emph{\Cref{sec:ethics} states the approval and the substance.
This appendix is the analysis behind it, organised by the department's six
domains and the Engineering Council's four pillars (honesty and integrity;
respect for life, law, the environment and the public good; accuracy and
rigour; leadership and communication), and it names every stakeholder.
Source: \repo{docs/research/03\_ethics\_and\_safety.md}, the document the
approval application was adapted from.}

\section{Stakeholders}

The \textbf{wearer}: bears \SI{17}{\kilo\gram} and the physical risk of
motion they do not command; the party with the least control and the most
at stake. The \textbf{operator}: bears the cognitive load and the
responsibility for motion they cannot feel. The \textbf{experimenter}: calls
breaks and stops, and must not let the schedule pressure either
participant. \textbf{Fellow students and the laboratory}: share the arms,
the room and the reputation of the group. \textbf{The College}: holds the
approval, the data and the liability. \textbf{The wider public}: the
eventual users of assistive limbs, and the people who will stand next to
them. \textbf{The environment}: the hardware's energy and material cost.

\section{A. Scientific integrity -- is it true?}

Every quantitative claim in the main body carries its evidence class
(implemented, verified in simulation, measured on hardware, or measured in
the pilot), and the report states that nothing below the pilot was
validated on a loaded physical arm. Twenty-six occasions on which the
measuring tool, not the system, was wrong are catalogued in
\cref{ch:instrument} and summarised in \cref{sec:validity}; each is
reported as a defect found, not hidden. Numbers that disagreed with their
source when this report was checked -- the e-stop count, the estimator's
wrong-guess rate at one cube pitch, the pinned-wrist reach ratio -- were
corrected and the correction is recorded in the repository log. The
analysis for the full study was pre-registered and run on \num{20000}
constructed data sets before any participant data existed
(\cref{app:analysis}), so its false-positive rate is known
(\SI{5.0}{\percent} against a nominal \SI{5}{\percent}).

\section{B. Scientific collegiality -- is it fair?}

The platform is shared laboratory equipment. Sessions were booked, the
two-stack rule (never two control processes on one robot) protected other
users' measurements as much as this project's, and every session leaves a
manifest and a log. The report anonymises participants and credits prior
work by citation; where this project reproduces a published finding it
says so, and where it contradicts one -- the robot's path being longer
under assistance -- it gives the measured mechanism rather than an
appeal.

\section{C. Protection of human subjects}

\textbf{Approval and consent.} SETREC approval 7111986. Participants
received an information sheet at least twenty-four hours in advance
stating what the device is, that it moves near them, its weight if worn,
the session length, the right to withdraw, what is recorded and how to
withdraw data afterwards. Consent is taken separately for participation,
for video (optional), for retention of anonymised data beyond the project,
and for its use in publication. Withdrawal is possible at any time without
reason, and data can be withdrawn up to publication through the
participant code.

\textbf{The two-person rule.} A wearer and an operator who arrive together
will not decline in front of each other, and the wearer consents to strictly
more. Consent is therefore taken in separate rooms, on separate forms, and
neither is told whether the other consented until both have. The wearer
holds their own emergency stop on the watchdog path and presses it twice,
watching the arms halt, before any recorded trial. No justification is ever
requested for a stop. A session ended by a wearer is reported as data, not
replaced with a fresh pair to make up the numbers, because that would select
for tolerant wearers.

\textbf{Risk register (the pilot's operator-only sessions, and the planned
worn study).} Robot contacts the participant: collision-aware solving, a
\SI{150}{\milli\metre} clearance floor enforced from the robot's true pose,
the autonomy never permitted to move toward the wearer, the wearer's own
stop. Unexpected or fast motion: a participant speed cap stricter than the
developer's, re-asserted every two seconds so a stray parameter change
cannot raise it; motion must be explicitly armed. Stop out of reach: the
software refuses to arm until the reach check is confirmed. A sensor fails
mid-trial: the trial is aborted and marked invalid, never silently
continued, because a frozen command looks like a slow participant in the
data. Physical strain from the master: blocks of at most eight minutes,
two-minute rests, a session of at most forty-five minutes, a Borg fatigue
probe between blocks. Strain from wearing the rig: twenty minutes maximum,
two people to don and doff, exclusion for any back, neck or shoulder
problem. Psychological risk: participants are told the \emph{system} is
under evaluation and that cancelling an autonomous action is normal use.

\textbf{What is not yet safe, stated rather than glossed.} The arms'
proximal links interfere with the wearer's own upper arms in the collision
model; this is set by the mount position and no planner can avoid it.
Worn operation therefore was not sought for approval; the pilot ran with a
mannequin wearer. The clearance check that runs on the physical arm samples
link origins, not link geometry (\cref{ch:discussion}) -- the one item
this report classes as a safety defect rather than an incompleteness.

\section{D. Animal welfare}

No animals are involved.

\section{E. Institutional integrity -- is it wise?}

Data handling: participant codes only; no name, email, date of birth,
address or telephone number appears in any file, enforced by
\repo{write\_manifest()} raising if one is passed; demographics as age
bands. The code-to-person mapping exists only to allow withdrawal, is held
on encrypted College storage outside the repository, and is destroyed at
publication. Video, if consented to, frames the rig and hands, not the
face. Anonymised data are retained under institutional policy; the
repository is not a data store. Legal basis: public-task processing under
UK GDPR and the Data Protection Act 2018, with consent for the optional
elements. No external sponsor or conflict of interest exists; the arms are
laboratory property and the findings belong to the College.

\section{F. Social responsibility}

The device this work points toward is an assistive limb worn by one person
and, in this configuration, directed by another. The result that
assistance shortens the operator's task while lengthening the robot's
path, with the wearer's trust and agency unmeasured, is exactly the shape
of harm that professional ethics exists to notice: a benefit to the party
who chooses and a cost to the party who does not. That is why H3 and H4 --
that the wearer's trust and agency fall under assistance -- are the
hypotheses the full study exists to test, and why the report refuses to
call the pilot evidence about the wearer. Sustainability and cost: the
platform re-uses existing laboratory arms and a consumer headset; the
\SI{4}{\gibi\byte} video-memory limit was met by choosing lighter methods
rather than buying hardware; the software is open in the laboratory
repository so the work can be continued rather than repeated.
